% !PRINTED 437
\weberpage{437}
%<<CONT
und $\alpha$, $\beta$, $\gamma$ sind die Wurzeln der cubischen Gleichung
\[ \alpha^3 + \alpha^2 - \frac{1}{7} = 0. \]
Man kann nun nach \S.\,72, (15), (12), (13)
\weq{15}{\Theta^2 = \tau^4\,\omega\,\tau^5\,\chi}
setzen, und daraus erhält man leicht nach (4) und (5)
\weq{16}{\Theta^2 = \begin{pmatrix} \beta, & \gamma\varepsilon^3, & \alpha\varepsilon^2 \\ \gamma\varepsilon^{-3}, & \alpha, & \beta\varepsilon^{-1} \\ \alpha\varepsilon^{-2}, & \beta\varepsilon, & \gamma \end{pmatrix}.}

Vergleicht man dies mit dem Resultate, was man erhält, wenn man aus (10) direct $\Theta^2$ bildet, so ergiebt sich:
\weq{17}{\begin{aligned}
\alpha &= \alpha^2\varepsilon^{-1} + \beta^2\varepsilon + \gamma^2\varepsilon^{-2}, & \alpha &= \beta\gamma + \gamma\alpha\varepsilon^2 + \alpha\beta\varepsilon^{-1} \\
\beta &= \alpha^2\varepsilon^{-1} + \beta^2\varepsilon^3 + \gamma^2\varepsilon^{-3}, & \beta &= \beta\gamma\varepsilon^3 + \gamma\alpha + \alpha\beta\varepsilon \\
\gamma &= \alpha^2\varepsilon^2 + \beta^2\varepsilon^3 + \gamma^2\varepsilon^{-2}, & \gamma &= \beta\gamma\varepsilon^{-3} + \gamma\alpha\varepsilon^{-2} + \alpha\beta,
\end{aligned}}

und die Vergleichung mit dem aus (10) gebildeten $\Theta^{-2}$ ergiebt ein ganz ähnliches Formelsystem, das aus (17) hervorgeht, wenn $\varepsilon$ mit $\varepsilon^{-1}$ vertauscht wird, nämlich:
\[ \begin{aligned}
\alpha &= \alpha^2\varepsilon + \beta^2\varepsilon^{-1} + \gamma^2\varepsilon^2, & \alpha &= \beta\gamma + \gamma\alpha\varepsilon^{-2} + \alpha\beta\varepsilon \\
\beta &= \alpha^2\varepsilon + \beta^2\varepsilon^{-3} + \gamma^2\varepsilon^3, & \beta &= \beta\gamma\varepsilon^{-3} + \gamma\alpha + \alpha\beta\varepsilon^{-1} \\
\gamma &= \alpha^2\varepsilon^{-2} + \beta^2\varepsilon^{-3} + \gamma^2\varepsilon^2, & \gamma &= \beta\gamma\varepsilon^3 + \gamma\alpha\varepsilon^2 + \alpha\beta.
\end{aligned} \]

Damit aber hierin kein Cirkelschluss liege, ist zu zeigen, dass die Relationen (17) in Folge der Bestimmungen (11), (12) wirklich erfüllt sind, denn dann erst ist das Bestehen der Relation (15) und die Gleichheit von $\Theta^2$ mit $\Theta^{-2}$ nachgewiesen. Es genügt dazu, zwei dieser Relationen, etwa
\[ \begin{aligned}
\alpha &= \alpha^2\varepsilon^{-1} + \beta^2\varepsilon + \gamma^2\varepsilon^{-2} \\
\alpha &= \beta\gamma + \gamma\alpha\varepsilon^2 + \alpha\beta\varepsilon^{-1},
\end{aligned} \]
abzuleiten; denn hat man diese nachgewiesen, so kann man darin $\varepsilon$ mit $\varepsilon^{-1}$ vertauschen, wodurch $\alpha$, $\beta$, $\gamma$ nicht geändert werden, und sodann $\varepsilon$ mit $\varepsilon^2$ und $\varepsilon^4$, wodurch $\alpha$, $\beta$, $\gamma$ in $\gamma$, $\alpha$, $\beta$ und in $\beta$, $\gamma$, $\alpha$ übergehen; und daher erhält man das ganze Formelsystem (17).

Diese beiden Relationen ergeben sich aber durch eine einfache Rechnung nach (11), (12) in der Form:
\[ \begin{aligned}
&-(\varepsilon + \varepsilon^2 + \varepsilon^4 - \varepsilon^{-1} - \varepsilon^{-2} - \varepsilon^{-4})(\varepsilon^4 - \varepsilon^{-4}) \\
&= \varepsilon^{-1}(\varepsilon^4 - \varepsilon^{-4})^2 + \varepsilon(\varepsilon^2 - \varepsilon^{-2})^2 + \varepsilon^{-2}(\varepsilon - \varepsilon^{-1})^2 \\
&= (\varepsilon - \varepsilon^{-1})(\varepsilon^2 - \varepsilon^{-2}) + \varepsilon^2(\varepsilon - \varepsilon^{-1})(\varepsilon^4 - \varepsilon^{-4}) \\
&\phantom{=} + \varepsilon^{-1}(\varepsilon^2 - \varepsilon^{-2})(\varepsilon^4 - \varepsilon^{-4}).
\end{aligned} \]
